%%% FORMATO E IDIOMA %%%
    \documentclass[letterpaper,DIV=15,12pt]{scrartcl}
    \usepackage[spanish,mexico,shorthands=off,es-lcroman]{babel}
    \usepackage{scrlayer-scrpage}
        \clearpairofpagestyles
        \ihead{\footnotesize \textit{Teoría de Conjuntos III}}
        \ohead{\footnotesize \textit{2026-2}}
        \cfoot{\normalfont\thepage}
        \addtokomafont{title}{\bfseries \rmfamily}
        \setlength{\headsep}{5pt}
        \newpairofpagestyles{beginstyle}{
            \clearpairofpagestyles
            \KOMAoptions{headsepline=false}
            \cfoot{\footnotesize \pagemark}
            \cfoot{\normalfont\thepage}
        }
        \addtokomafont{section}{\raggedleft\rmfamily}
        \addtokomafont{subsection}{\raggedleft\rmfamily}
        \addtokomafont{subsubsection}{\raggedleft\rmfamily}
        \setlength{\headsep}{8pt}
        \setlength{\footskip}{20pt}
        \setlength{\textheight}{600pt}
%%% PAQUETES %%%
    \usepackage{ifthen}
    \usepackage[shortlabels]{enumitem}
        \setenumerate[1]{label=\MakeLowercase{\roman*}), ref=\roman*}
        \setenumerate[2]{label=\MakeLowercase{\alph*}), ref=\alph*}
    \usepackage{amsmath}
    \usepackage{amsthm}\usepackage[T1]{fontenc}
    \usepackage{stix2}
	\renewcommand{\qedsymbol}{$\blacksquare$}
	\addto\captionsspanish{\renewcommand{\proofname}{\normalfont\bfseries Demostración}}	
	\newenvironment{solucion}{\renewcommand{\qedsymbol}{$\boxtimes$}\begin{proof}[\normalfont\bfseries Solución]}{\end{proof}}

    \usepackage{amsthm}
		\newtheoremstyle{manual}{3pt}{3pt}{\normalfont}{}{\bfseries}{.}{.5em}{#1 #3}
		\theoremstyle{manual}
		\newtheorem{proposicion}{Proposición}
        \newtheorem{definicion}{Definición}
        \newtheorem{teorema}{Teorema}
        \newtheorem{lema}{Lema}
        \newtheorem{corolario}{Corolario}
        \newtheorem*{ejercicio*}{Ejercicio}        
%%% C O M A N D O S %%%
    \renewcommand{\emptyset}{\varnothing}
    \newcommand{\tq}{:}
%% D O C U M E N T O %%
\begin{document}
\thispagestyle{plain}

    \begin{center}
        {\fontsize{30}{60}\rmfamily \textbf{Más sobre modelos ET}} \\ \vspace{.2cm} Teoría de Conjuntos III, 2026-2
    \end{center}
    \begin{flushright}
        \footnotesize \hfill Profesor: Luis Jesús Turcio Cuevas.\\
        \hfill Ayudante: Hugo Víctor García Martínez.
    \end{flushright}

    A continuación se prueban los siguientes criterios que permiten determinar cuándo un modelo transitivo estándar satisface los axiomas de potencia y separación*. Utilizaremos que las nociones de ``ser elemento, subconjunto, ser vacío, etc...'' son $M$-$V$ absolutas para clases transitivas $M$.
    
    \begin{proposicion}[1]\itshape
        Sea $M$ una calse (no necesariamente propia) transitiva.
        \begin{enumerate}
            \item Si para cada $x \in M$ y cada $y \subseteq x$, $y \in M$, \textbf{entonces} $M$ modela el esquema de separación.
            \item $M$ modela el axioma de separación \textbf{si y sólo si} para cada $x \in M$ ocurre $\mathscr{P}(x) \cap M \in M$.
        \end{enumerate}
    \end{proposicion}
    \begin{proof}\textbf{(i)}
        Supóngase que, si $x \in M$ y $y \subseteq x$, entonces $y \in M$; es decir, $\mathscr{P}(x) \subseteq M$. Sea $\varphi$ una fórmula de la teoría de conjuntos, veamos que
        \[ M \vDash \forall x \, \exists y \, \forall w \, \Big( w \in y \leftrightarrow \big( w \in x \land \varphi(w) \big) \Big) \, , \]
        es decir, que ``$\forall x \in M \, \exists y \in M \, \forall w \, (w \in y \leftrightarrow (w \in x \land \varphi^M(w)))$'' se cumple.

        En efecto, sean $x \in M$ y $y:=\{ w \in x \tq \varphi^M (w) \}$, entonces $y$ es conjunto (en $V$) por separación (de fuera). Por hipótesis, $y \in M$, claramente, para cada $w \in M$ ocurre que $w \in y$ si y sólo si $w \in x$ y $\varphi^M(w)$. Por lo tanto $M \vDash \operatorname*{\textsc{Sep}}^*$.
    \end{proof}
    \begin{proof}\textbf{(ii)}
        Supóngase primero que $M \vDash \operatorname*{\textsc{Pot}}$ y sea $x \in M$, entonces existe $y \in M$ tal que para cada $w \in M$, $w \in y$ ocurre si y sólo si $(w \subseteq x)^M \equiv w \subseteq x$. Como $M$ es transitiva, $y \subseteq M$, entonces de lo anterior, para todo conjunto $w$ (en $V$), $w \in y$ si y sólo si $w \in M$ y $w \subseteq x$, esto es, $y = \mathscr{P}(x) \cap M$.

        ($\leftarrow$) Ahora, supóngase que para cada $x \in M$ ocurre $\mathscr{P}(x) \cap M \in M$, veamos que:
        \[ M \vDash \forall x \, \exists y \, \forall w \, \big( w \in y \leftrightarrow ( w \subseteq x ) \big) \, , \]
        es decir, que ``$\forall x \in M \, \exists y \in M \, \forall w \in M \, \big( w \in y \leftrightarrow ( w \subseteq x ) \big)$'' es verdadera.

        Efectivamente, sean $x \in M$ y $y=\mathscr{P} \cap M$. Por hipótesis, $y \in M$, además, para cada $w \in M$ ocurre que $w \in y$ si y sólo si $w \in M$ (lo cual ya ocurre) y $x \in \mathscr{P}(x)$, es decir, $w \subseteq x$. Por lo tanto, se concluye que $M \vDash \operatorname*{\textsc{Pot}}$.
    \end{proof}

    El recíproco del punto (i) de la proposión anterior es falso en general. Se ha adelantado que existe un modelo estándar, transitivo y numerable $M$ de (buena parte de) \textsc{Zfc}, incluyendo todo el esquema de separación* y el axioma de infinito. Este modelo, por tanto, tiene a $\omega$ como elemento (recuerde que la fórmula $x=\omega$ es $M$-$V$ absoluta), por ello, no puede ser que $\mathscr{P}(\omega) \subseteq M$, esto contradiría la numerabilidad de $M$.

    Conviene entonces reflexionar sobre al menos dos cosas, dado lo comentado en el párrafo anterior, la primera es notar que los métodos para demostrar que un modelo, estándar y transitivo, satisface el axioma de separación*, pueden variar bastante, dependiendo del modelo. La segunda es preguntarse para qué tan útil resulta ese punto (i). Exploremos esto último.    

    Junto con los criterios vistos en clase, se puede utilizar la proposición anterior para concluir que $\operatorname*{BF}$ modela todo \textsc{Zfc}; y también, que si $\gamma$ es ordinal límite, entonces $\operatorname*{BF}_\gamma \vDash \operatorname*{\textsc{Zfc-Rmp}}^*$. Encontremos un modelo (conjunto, en $V$) más, que nos permite dar otro resultado de (no) equiconsistencia.

    \begin{definicion}[2]\itshape
        La \textbf{cerradura transitiva} de un conjunto $x$, denotada $\operatorname*{tcl}(x)$, es el $\subseteq$-mínimo conjunto transitivo $T$ tal que $x \subseteq T$.

        Este operador es claramente monótono, e idempotente para conjuntos transitivos.
    \end{definicion}
    
    Dentro de \textsc{Zfc}, se puede mostrar que la cerradura transitiva existe para cualquier conjunto. Es sencillo además, demostrar que para cada conjunto $x$:
    \begin{equation}\label{ident}
        \operatorname*{tcl}(x) = x \cup \bigcup_{y \in x} \operatorname*{tcl}(y)
    \end{equation}
    Nótese que si $w \in x \cup \bigcup_{y \in x} \operatorname*{tcl}(y)$, entonces $w \in x \subseteq \operatorname*{tcl}(x)$, o bien, existe $y \in x$ para el cual $w \in \operatorname*{tcl}(y)$. En este último caso, nótese que $y \in \operatorname*{tcl}(x)$, así que $y \subseteq \operatorname*{tcl}(x)$, y consecuentemente, $\operatorname*{tcl}(y) \subseteq \operatorname*{tcl}(\operatorname*{tcl}(x)) = \operatorname*{tcl}(x)$. Mostrando que, en cualquier caso $w \in \operatorname*{tcl}(x)$, por lo que ocurre la contención $\supseteq$ de (\ref{ident}). Para la contención recíproca, basta ver que el conjunto $x \cup \bigcup_{y \in x} \operatorname*{tcl}(y)$ es transitivo (ya que, por definición, contiene a $x$), esto se deja moral :)
    

    \begin{definicion}[3]\itshape
        Para cada cardinal $\kappa$, se define $\operatorname*{H}(\kappa)=\{ x \tq |\operatorname*{tcl}(x)| < \kappa  \}$. Particularmente, $\operatorname*{H}(\aleph_1)$ es el conjunto de conjuntos \textbf{hereditariamente contables}, se denota por $\operatorname*{HC}$.
    \end{definicion}

    Bajo el axioma de buena buena fundación, cada $H(\kappa)$ es conjunto, pues se separa (axioma de separación*) del conjunto $\operatorname*{BF}_\kappa$.

    \begin{lema}[4]
        El esquema de reemplazo, implica el de separación.
    \end{lema}
    \textit{\bfseries Hint.} Dada una fórmula $\varphi(x)$, se define el funcional $\psi(x,y) \leftrightharpoons (x=y \land \varphi(x))$. Para cada conjunto $A$, la imagen de $A$ bajo tal funcional, es $\{x \in A \tq \varphi(x)\}$. \hfill $\boxtimes$.

    El lema anterior facilita la prueba de lo siguiente.

    \begin{proposicion}[5]
        (Bajo \textsc{Zfc}) $\operatorname*{HC}$ es un modelo (conjunto) de $\operatorname*{\textsc{Zfc-Pot}}$.
    \end{proposicion}
    \begin{proof}
        Nótese que $\operatorname*{HC}$ es un conjunto transitivo, pues, para cada $x \in \operatorname*{HC}$ y cada $y \in x$, se tiene que $y \subseteq X$. Así, $\operatorname*{tcl}(y) \subseteq \operatorname*{tcl}(x)$, y con ello, que $y \in \operatorname*{HC}$.

        Los axiomas de extensionalidad, par, unión, infinito, son sencillos de verificar. Nótese que vacío se satisfará, pues hay axioma del infinito.
        
        Para el axioma de elección, sea $x \in \operatorname*{HC}$ un conjunto no vacío, de conjuntos no vacíos, entonces por \textsc{AC} (en $V$), $x$ admite un conjunto de elección $A \subseteq x$ (noción absoluta para modelos transitivos). Como $A \subseteq x$, entonces $\operatorname*{tcl}(A) \subseteq \operatorname*{tcl}(x)$; y con ello, $A \in \operatorname*{HC}$.

        Restan únicamente, los esquemas de separación* y reemplazo*. Supóngase que $f$ una función tal que $x=\operatorname*{dom}(f) \in \operatorname*{HC}$ y supóngase que $\operatorname*{im}(f) \subseteq \operatorname*{HC}$, por un teorema visto en clase, basta verificar que $\operatorname*{im}(f) \in \operatorname*{HC}$.

        Y en efecto, de (\ref{ident}) se obtiene que:
        \[ \operatorname*{tcl}(\operatorname*{im}(f)) = \operatorname*{im}(f) \cup \bigcup_{y \in x} \operatorname*{tcl}(f(y)) \]

        Nótese que, $x \subseteq \operatorname*{tcl}(x) \preccurlyeq \omega$, consecuentemente, $\operatorname*{im}(f)$ es contable. Ahora, dado $y \in x$, se tiene que $f(y) \in \operatorname*{im}(f) \subseteq \operatorname*{HC}$, así que $\operatorname*{tcl}(f(y)) \preccurlyeq \omega$. Así, de la igualdad anterior se tiene que $\operatorname*{tcl}(\operatorname*{im}(f))$ es un conjunto contable, unido con, una unión numerable de conjuntos contables. Por lo tanto $|\operatorname*{tcl}(\operatorname*{im}(f))| \leq \omega$, mostrando que $\operatorname*{im}(f) \in \operatorname*{HC}$. Así, $\operatorname*{HC}$ modela el esquema de reemplazo* (así mismo, el de separación*).
    \end{proof}

    Obsérvese que en la prueba anterior, lo único ``especial'' de $\aleph_1$ que se utilizó, fue su regularidad (unión contable de conjuntos contables es contable). Así que, en realidad, cualquier cardinal regular $\kappa$ tiene la propiedad de que $\operatorname*{H}(\kappa) \vDash \operatorname*{\textsc{Zfc - Rmp}^*}$.

    También, nótese que de no tener \textsc{AC} en $V$, se sigue pudiendo demostrar que $\operatorname*{HC}$ es modelo de todo $\operatorname*{\textsc{Zf}-Rmp}^*$. Por lo que el siguiente corolario vale también, escribiendo \textsc{Zf} en vez de \textsc{Zfc}.

    \begin{corolario}[6]
        Se cumple lo siguiente, bajo \textsc{Zfc}:
        \begin{enumerate}
            \item $\operatorname*{HC} \vDash \operatorname*{\textsc{Zfc-Pot}}$ y $\operatorname*{HC} \vDash \lnot \operatorname*{\textsc{Pot}}$.
            \item $\operatorname*{\textsc{Zfc-Pot}} \not\vdash \operatorname*{\textsc{Pot}}$, es decir, potencia ``sí era un axioma''.
            \item $\operatorname*{\textsc{Zfc-Pot}} \lhd \operatorname*{\textsc{Zfc}}$, es decir, potencia es un axioma que ``requiere más credibilidad''.
        \end{enumerate}
    \end{corolario}
    \begin{proof}
        Lo único que hay que probar es (i), lo haremos por ``golazo''. Si $\operatorname*{HC} \not\vDash \lnot \operatorname*{\textsc{Pot}}$, entonces $\operatorname*{HC} \vDash \operatorname*{\textsc{Pot}}$ (pues potencia es enunciado). Entonces $\operatorname*{HC} \vDash \operatorname*{\textsc{Zfc}}$; y como $\operatorname*{HC}$ se contruye dentro de $\operatorname*{\textsc{Zfc}}$, esto indica que en todo modelo de $\operatorname*{\textsc{Zfc}}$, existe un modelo (conjunto) de $\operatorname*{\textsc{Zfc}}$. Por lo tanto $\operatorname*{\textsc{Zfc}} \vdash \operatorname*{Con}(\operatorname*{\textsc{Zfc}})$, lo cual contradice el segundo teorema de incompletud de Gödel.
    \end{proof}

\end{document}                                                                                                     